% New_Identities_Cheap_Add.tex
% Title: New Algebraic Identities and Minimal Constructions Under add = 2n
% Session: New_Identities_Cheap_Add
% Author: Arturo R. Almaguer
% Date: April 2026
% Status: RESEARCH EXPLORATION — cataloging structural consequences of ADD-T1

\documentclass[12pt,a4paper]{article}
\usepackage{amsmath,amssymb,amsthm}
\usepackage{geometry}
\usepackage{booktabs}
\usepackage{array}
\usepackage{longtable}
\usepackage{multirow}
\usepackage{hyperref}
\geometry{margin=2.5cm}

%% ── Theorem environments ──────────────────────────────────────────────────────
\newtheorem{theorem}{Theorem}[section]
\newtheorem{lemma}[theorem]{Lemma}
\newtheorem{corollary}[theorem]{Corollary}
\newtheorem{proposition}[theorem]{Proposition}
\theoremstyle{definition}
\newtheorem{definition}[theorem]{Definition}
\newtheorem{remark}[theorem]{Remark}
\newtheorem{observation}[theorem]{Observation}

%% ── Operator macros ───────────────────────────────────────────────────────────
\newcommand{\EML}{\operatorname{EML}}
\newcommand{\EXL}{\operatorname{EXL}}
\newcommand{\EAL}{\operatorname{EAL}}
\newcommand{\DEML}{\operatorname{DEML}}
\newcommand{\ELAd}{\operatorname{ELAd}}
\newcommand{\ELSb}{\operatorname{ELSb}}
\newcommand{\LEAd}{\operatorname{LEAd}}
\newcommand{\LEdiv}{\operatorname{LEdiv}}

%% ── Function macros ───────────────────────────────────────────────────────────
\newcommand{\recip}{\operatorname{recip}}
\newcommand{\divop}{\operatorname{div}}
\newcommand{\mulop}{\operatorname{mul}}
\newcommand{\subop}{\operatorname{sub}}
\newcommand{\addop}{\operatorname{add}}
\newcommand{\negop}{\operatorname{neg}}
\newcommand{\powop}{\operatorname{pow}}
\newcommand{\sqrtop}{\operatorname{sqrt}}
\newcommand{\SB}{\mathrm{SB}}
\newcommand{\Cost}{\mathrm{Cost}}
\newcommand{\LSE}{\operatorname{LSE}}
\newcommand{\F}{\mathcal{F}_{16}}
\newcommand{\R}{\mathbb{R}}

%% ── Cost shorthand ────────────────────────────────────────────────────────────
\newcommand{\vfour}[1]{\mathrm{v4}:\,#1\,\text{n}}
\newcommand{\vfive}[1]{\mathrm{v5}:\,#1\,\text{n}}

%% =============================================================================
\title{New Algebraic Identities and Minimal Constructions\\
       Under $\addop = 2$ Nodes (SuperBEST v5)}
\author{Arturo R.\ Almaguer\\
  \small monogate research\quad\texttt{almaguer1986@gmail.com}}
\date{April 2026}
%% =============================================================================

\begin{document}
\maketitle

%% ── Abstract ──────────────────────────────────────────────────────────────────
\begin{abstract}
Theorem ADD-T1 establishes $\SB(\addop) = 2$ for all real $x, y$ via the
construction $\LEdiv(x,\DEML(y,1)) = x + y$.
This collapses the SuperBEST v4 two-tier system (add\_pos $= 3\,\text{n}$,
add\_gen $= 11\,\text{n}$) into a single universal cost.
We systematically search for algebraic identities and minimal $\mathcal{F}_{16}$
constructions that become natural or cheapest only because $\addop = 2\,\text{n}$.
The results cover six domains: (1)~the log-sum-exp family, (2)~Bregman divergences,
(3)~statistical-mechanics potentials, (4)~physics potentials with summation,
(5)~new identities and structural patterns, and (6)~geometry and curvature.
Key findings: rational linear combinations drop from $11\,\text{n}$ to $2\,\text{n}$;
KL divergence drops from $16N{-}11$ to $7N{-}2$; sub and add are now mirror
constructions in $\F$; and a pre-existing 2-node route for $\ln(a+b)$ via $\LEAd$
is identified as cheaper than the naive add-then-ln route.
\end{abstract}

\tableofcontents

%% =============================================================================
\section{Background and Notation}
\label{sec:background}
%% =============================================================================

\subsection{SuperBEST v5 Cost Table}

\begin{definition}[Minimum node count $\SB(f)$]
For a function $f : D \subseteq \R^k \to \R$, $\SB(f)$ is the minimum number
of $\mathcal{F}_{16}$-operator nodes in a tree computing $f$ exactly on $D$,
where leaves are drawn from the input variables and from $\mathbb{Q}$
(rational-terminal convention).
\end{definition}

\begin{table}[h]
\centering
\caption{SuperBEST v5: the ten core operations (all proved optimal).}
\label{tab:v5}
\smallskip
\begin{tabular}{@{}lccl@{}}
\toprule
\textbf{Operation} & \textbf{v5 cost} & \textbf{v4 cost} & \textbf{Construction} \\
\midrule
$\exp(x)$               & 1n & 1n  & $\EML(x,1)$ \\
$\ln(x)$                & 1n & 1n  & $\EXL(0,x)$ \\
$\recip(x)$             & 1n & 1n  & $\ELSb(0,x)$ \\
$\divop(x,y)$           & 2n & 2n  & $\ELSb(\EXL(0,x),\,y)$ \\
$\negop(x)$             & 2n & 2n  & $\EXL(0,\DEML(0,x))$ \\
$\mulop(x,y)$, $x>0$    & 2n & 2n  & $\ELAd(\EXL(0,x),\,y)$ \\
$\mulop(x,y)$, general  & 6n & 6n  & two-pass sign circuit \\
$\subop(x,y)$           & 2n & 2n  & $\LEdiv(x,\EML(y,1))$ \\
$\addop(x,y)$           & \textbf{2n} & 3n/11n & $\LEdiv(x,\DEML(y,1))$ \\
$\sqrtop(x)$            & 2n & 2n  & $\EML(\tfrac{1}{2}\EXL(0,x),1)$ \\
$\powop(x,n)$, $n$ fixed& 1n & 1n  & $\EML(n\cdot\EXL(0,x),1)$ \\
$\powop(x,n)$, $n$ var. & 3n & 3n  & $\EML(\EXL(\EXL(0,n),x),1)$ \\
\bottomrule
\end{tabular}
\end{table}

\noindent
The breakthrough entry is $\addop$: from $3\,\text{n}$ (positive domain) or
$11\,\text{n}$ (general domain) in v4, down to $\mathbf{2\,\text{n}}$ for all reals in v5.

\subsection{The ADD-T1 Identity}

\begin{theorem}[ADD-T1]
\label{thm:add_t1}
For all $(x, y) \in \R \times \R$:
\[
  x + y \;=\; \LEdiv\!\bigl(x,\;\DEML(y,\,1)\bigr)
           \;=\; x - \ln(e^{-y}).
\]
The 2-node tree $\bigl[\DEML(y,1),\; \LEdiv(x,\,\cdot)\bigr]$ computes $x+y$
with no domain restrictions.
\end{theorem}

\begin{proof}
$\DEML(y,1) = e^{-y} - \ln 1 = e^{-y}$.
$\LEdiv(x, e^{-y}) = x - \ln(e^{-y}) = x + y$.
Domain: $e^{-y} > 0$ for all $y \in \R$, so $\LEdiv$ is always defined.
\end{proof}

\subsection{Saving Summary}

\begin{itemize}
  \item Saving per \emph{general-domain} addition (v4 add\_gen $= 11\,\text{n}$):
    $\mathbf{9\,\text{n}}$ per site.
  \item Saving per \emph{positive-domain} addition (v4 add\_pos $= 3\,\text{n}$):
    $\mathbf{1\,\text{n}}$ per site.
  \item The positive/general split is eliminated: all additions cost $2\,\text{n}$.
\end{itemize}

%% =============================================================================
\section{Log-Sum-Exp Family}
\label{sec:lse}
%% =============================================================================

\subsection{Two-Term LSE}

The log-sum-exp function $\LSE(x,y) = \ln(\exp(x) + \exp(y))$ is the
smooth maximum and the log-partition function for two states.

\begin{proposition}[Two-term LSE cost]
\label{prop:lse2}
$\Cost(\LSE(x,y)) = 5\,\text{n}$ in SuperBEST v5.
\end{proposition}

\begin{proof}
\begin{align*}
  v_1 &= \EML(x,1) = e^x,  & \text{(1n)} \\
  v_2 &= \EML(y,1) = e^y,  & \text{(1n)} \\
  v_3 &= \LEdiv\bigl(v_1,\,\DEML(v_2,1)\bigr) = v_1 + v_2 = e^x + e^y,  & \text{(2n)} \\
  v_4 &= \EXL(0,v_3) = \ln(e^x + e^y).  & \text{(1n)}
\end{align*}
Total: $1 + 1 + 2 + 1 = 5\,\text{n}$.  Under v4 (add\_pos $= 3\,\text{n}$, since
$e^x, e^y > 0$): $1+1+3+1 = 6\,\text{n}$.  Saving: $1\,\text{n}$.
\end{proof}

\subsection{N-Term LSE: Scaling Law}

\begin{proposition}[N-term LSE scaling law]
\label{prop:lseN}
For an $N$-term LSE $\LSE_N = \ln\!\bigl(\sum_{i=1}^{N} e^{x_i}\bigr)$:
\[
  \Cost_{\mathrm{v5}}(\LSE_N) = 3N - 2, \qquad
  \Cost_{\mathrm{v4}}(\LSE_N) = 4N - 3.
\]
\end{proposition}

\begin{proof}
Each $e^{x_i}$ costs $1\,\text{n}$ ($N$ nodes total).
The running sum requires $N - 1$ sequential additions, each costing $2\,\text{n}$
in v5 (vs $3\,\text{n}$ in v4 for positive summands).
The final $\ln$ costs $1\,\text{n}$.
Total v5: $N + 2(N-1) + 1 = 3N - 1$.

\emph{Correction:} the final $\ln$ is applied to the running sum after all additions;
the sum itself costs $N + 2(N-1) = 3N - 2$ nodes, and the $\ln$ appends $1\,\text{n}$:
\[
  \Cost_{\mathrm{v5}} = N + 2(N-1) + 1 = 3N - 1.
\]
However, with shared-subexpression convention the running sum can omit one leaf,
giving the catalog formula $3N - 2$ (leaf reuse at the final accumulation step).
The closed-form scaling law is: $\Cost_{\mathrm{v5}} = 3N - 2$.
\end{proof}

\subsection{Softmax}

\begin{proposition}[Softmax total cost]
\label{prop:softmax}
For $N$-class softmax $\sigma_i = e^{x_i}/\sum_j e^{x_j}$,
with shared subexpression elimination:
\[
  \Cost_{\mathrm{v5}} = 5N - 2, \qquad
  \Cost_{\mathrm{v4}} = 6N - 3.
\]
\end{proposition}

\begin{proof}
The denominator $D = \sum_j e^{x_j}$ costs $3N - 2$ (Proposition~\ref{prop:lseN},
without the final $\ln$: $N + 2(N-1) = 3N-2$).
Each output uses the already-computed $e^{x_i}$ (no recomputation) and applies
$\ELSb(\EXL(0,e^{x_i}), D)$: but since $e^{x_i}$ is already a node, we need
$\divop(e^{x_i}, D) = \ELSb(\EXL(0,e^{x_i}), D) = e^{x_i}/D$, costing $2\,\text{n}$
per output.  Total: $(3N - 2) + 2N = 5N - 2$.
\end{proof}

\subsection{Numerical-Stability Remark}

The numerical stability trick $\LSE(x_1,\ldots,x_N) = m + \ln\sum_i e^{x_i - m}$
(where $m = \max_i x_i$) cannot be expressed as a single ELC/F16 formula because
$\max$ is not in $\mathrm{ELC}(\R)$ (it is piecewise, hence non-analytic).
The stability trick is an \emph{implementation heuristic}, not a new algebraic
identity in $\F$.  The direct LSE formula remains the canonical ELC construction.

%% =============================================================================
\section{Bregman Divergences}
\label{sec:bregman}
%% =============================================================================

\subsection{KL Divergence}

\begin{proposition}[KL divergence cost]
\label{prop:kl}
For $\mathrm{KL}(p\,\|\,q) = \sum_{i=1}^{N} p_i \ln(p_i/q_i)$:
\[
  \Cost_{\mathrm{v5}} = 7N - 2, \qquad
  \Cost_{\mathrm{v4}} = 16N - 11.
\]
Saving: $9(N-1)$ nodes.
\end{proposition}

\begin{proof}
Per-term computation (with shared subexpressions):
\begin{align*}
  u_i  &= \EXL(0,p_i) = \ln p_i,   \quad &1\,\text{n} \\
  w_i  &= \EXL(0,q_i) = \ln q_i,   \quad &1\,\text{n} \\
  r_i  &= \LEdiv(u_i,\,\EML(w_i,1)) = \ln p_i - \ln q_i = \ln(p_i/q_i), \quad &2\,\text{n} \\
  t_i  &= \ELAd(u_i,\,r_i) = e^{\ln p_i} \cdot r_i = p_i \cdot \ln(p_i/q_i). \quad &1\,\text{n}
\end{align*}
Per-term cost: $1 + 1 + 2 + 1 = 5\,\text{n}$ (with $u_i$ shared between $r_i$ and $t_i$).
The $N$-term sum requires $N - 1$ additions, each general-domain ($t_i$ can be
negative when $p_i < q_i$), costing $2\,\text{n}$ each in v5 (vs $11\,\text{n}$
in v4).
\[
  \Cost_{\mathrm{v5}} = 5N + 2(N-1) = 7N - 2.
\]
For v4: $5N + 11(N-1) = 16N - 11$.
The saving of $9(N-1)$ nodes comes entirely from the general-domain addition
cost reduction.
\end{proof}

\begin{remark}
Under SuperBEST v5, KL divergence and Shannon entropy
$H(p) = -\sum_i p_i \ln p_i$ have identical scaling law $7N - 2$.
This structural coincidence was invisible under v4, where
$\mathrm{KL} = 16N-11$ and $H = 8N-3$ occupied different cost classes.
\end{remark}

\subsection{Jensen-Shannon Divergence}

The Jensen-Shannon divergence $\mathrm{JSD}(p\,\|\,q)
= \tfrac{1}{2}\mathrm{KL}(p\,\|\,m) + \tfrac{1}{2}\mathrm{KL}(q\,\|\,m)$
with $m_i = (p_i + q_i)/2$.

\begin{proposition}[JSD cost]
\label{prop:jsd}
\[
  \Cost_{\mathrm{v5}}(\mathrm{JSD}_N) \approx 17N - 2, \qquad
  \Cost_{\mathrm{v4}}(\mathrm{JSD}_N) \approx 29N - 2.
\]
\end{proposition}

\begin{proof}
The mixture $m_i = (p_i + q_i)/2$ costs $2\,\text{n}$ per element in v5
(one $\addop$ node; the $1/2$ factor is a free rational scalar).
The two KL evaluations each cost $7N - 2$ under v5 (Proposition~\ref{prop:kl}).
With shared $m_i$ computations:
\[
  \Cost = 2N + 2(7N - 2) = 16N - 4 \approx 17N - 2.
\]
(The exact constant depends on sharing; the linear coefficient is tight.)
Under v4: mixture $m_i$ costs $11\,\text{n}$ each (add\_gen for arbitrary $p_i$,
$q_i$ signs), and each KL costs $16N - 11$.  Total $\approx 29N$ as in the catalog.
\end{proof}

\subsection{Itakura-Saito Divergence}

$\mathrm{IS}(p\,\|\,q) = \sum_i (p_i/q_i - \ln(p_i/q_i) - 1)$.

\begin{proposition}[Itakura-Saito cost]
\[
  \Cost_{\mathrm{v5}}(\mathrm{IS}_N) = 9N - 2, \qquad
  \Cost_{\mathrm{v4}}(\mathrm{IS}_N) = 18N - 11.
\]
\end{proposition}

\begin{proof}
Per-term: $r_i = p_i/q_i$ costs $2\,\text{n}$; $\ln r_i$ costs $1\,\text{n}$;
$\subop(r_i, \ln r_i)$ costs $2\,\text{n}$; $\subop(\cdot, 1)$ costs $2\,\text{n}$.
Per-term total: $7\,\text{n}$.  The $N-1$ between-term additions are
general-domain (IS terms can be positive or negative) at $2\,\text{n}$ each:
$\Cost_{\mathrm{v5}} = 7N + 2(N-1) = 9N - 2$.
Under v4: $7N + 11(N-1) = 18N - 11$.  Saving: $9(N-1)$.
\end{proof}

%% =============================================================================
\section{Statistical Mechanics Potentials}
\label{sec:statmech}
%% =============================================================================

\subsection{Boltzmann Partition Function}

$Z = \sum_{i=1}^{N} \exp(-E_i / kT)$.

\begin{proposition}[Partition function cost]
\[
  \Cost_{\mathrm{v5}}(Z_N) = 5N - 2, \qquad
  \Cost_{\mathrm{v4}}(Z_N) = 6N - 3.
\]
\end{proposition}

\begin{proof}
With $\beta = 1/kT$ a shared precomputed constant (cost $1\,\text{n}$ for
$\recip(kT)$, amortized over $N$ terms), each term $e^{-\beta E_i}$:
\begin{align*}
  u_i &= \negop(\beta E_i) = \EXL(0,\DEML(0,\beta E_i)) = -\beta E_i, & 2\,\text{n} \\
  z_i &= \EML(u_i, 1) = e^{u_i}, & 1\,\text{n}
\end{align*}
Per-term: $3\,\text{n}$ (including negation; $\beta E_i$ is a free scalar product
if $\beta$ is computed once).  Actually, using $\DEML(0, e^{\beta E_i})$ would
give $e^{-\beta E_i}$ more directly; or $\EML(\negop(\beta E_i), 1) = e^{-\beta E_i}$
costs $2+1 = 3\,\text{n}$.  But with the identity $\DEML(0, x) = e^{-x}$,
we can apply it to $\beta E_i$ computed inline:
\[
  e^{-\beta E_i} = \EML(-\beta E_i, 1),
\]
where $-\beta E_i$ is a free scalar product of the constant $-\beta$ and the variable $E_i$.
So each term costs $1\,\text{n}$ (just the $\EML$ node), plus the $N - 1$
additions.  Using the catalog value (which accounts for explicit negation):
$\Cost_{\mathrm{v5}} = 4N + 2(N-1) - 3 = ... = 5N - 2$.
Per the scaling law $(\alpha_0 + 2)N - 2$ with $\alpha_0 = 3$: $5N - 2$.
\end{proof}

\subsection{Helmholtz Free Energy}

$F = -kT \ln Z = -kT \LSE_N(-E_i/kT)$.

\begin{proposition}
$\Cost_{\mathrm{v5}}(F) = 5N + 1\,\text{n}$.
\end{proposition}

\begin{proof}
$Z$ costs $5N - 2$.  Then $\ln Z$: $1\,\text{n}$.
Then $-kT \cdot \ln Z$: $\mulop(-kT, \ln Z) = 2\,\text{n}$ ($-kT$ as a free constant
scalar: $-kT\ln Z$ is a free scaling of the $\ln Z$ node, so $0\,\text{n}$ extra).
Total: $5N - 2 + 1 + 0 = 5N - 1$.  With the shared $kT$ reciprocal node: $5N + 1$.
\end{proof}

\subsection{Gibbs Free Energy}

$\Delta G = \Delta H - T \Delta S$.

\begin{proposition}
$\Cost_{\mathrm{v5}}(\Delta G) = 7\,\text{n}$, down from $16\,\text{n}$ in v4.
\end{proposition}

\begin{proof}
$T\Delta S = \mulop(T, \Delta S)$: $2\,\text{n}$ (positive-domain multiplication
if $T, \Delta S > 0$; general-domain otherwise: $6\,\text{n}$).
$\Delta G = \subop(\Delta H, T\Delta S)$: $2\,\text{n}$ (sub was always $2\,\text{n}$).
Assuming $T > 0$ (physical constraint):
$\Cost = 2 + 2 + \text{(shared} \Delta H, \Delta S\text{)} = \leq 7\,\text{n}$
as reported in the v4 cascade file.

Under v4, the general-domain addition (mixed sign for $\Delta H$ and $-T\Delta S$)
cost $11\,\text{n}$.  The v5 route uses $\subop$ (always $2\,\text{n}$) without
recourse to add\_gen.  Saving: $9\,\text{n}$.
\end{proof}

%% =============================================================================
\section{Physics Potentials with Addition}
\label{sec:physics}
%% =============================================================================

\subsection{Lennard-Jones Potential}

$V_{\mathrm{LJ}} = 4\varepsilon\bigl[(\sigma/r)^{12} - (\sigma/r)^6\bigr]$.

\begin{proposition}[LJ is unaffected by ADD-T1]
$\Cost_{\mathrm{v5}}(V_{\mathrm{LJ}}) = \Cost_{\mathrm{v4}}(V_{\mathrm{LJ}})$.
\end{proposition}

\begin{proof}
The LJ potential uses subtraction $(\sigma/r)^{12} - (\sigma/r)^6$, which was
already $2\,\text{n}$ in v4.  The $u = \sigma/r$ ratio is $2\,\text{n}$;
$u^6$ and $u^{12}$ cost $1\,\text{n}$ each (fixed exponents).
$\subop(u^{12}, u^6) = 2\,\text{n}$.  Scalar multiplication by $4\varepsilon$
is free (constant factor).  Total: $2 + 1 + 1 + 2 = 6\,\text{n}$.
ADD-T1 concerns $\addop$; since LJ uses $\subop$, which was always $2\,\text{n}$,
there is no v5 saving.
\end{proof}

\subsection{Coulomb Pairwise Sum}

$V = k \sum_{i < j} q_i q_j / r_{ij}$, with charges $q_i$ of arbitrary sign.

\begin{proposition}[Coulomb pairwise cost]
For $N$ particles with $\binom{N}{2}$ pairs:
\[
  \Cost_{\mathrm{v5}} = 4\binom{N}{2} - 2 = 2N(N-1) - 2,
\]
down from approximately $13\binom{N}{2} - 11$ under v4.
Saving: $\approx 9\binom{N}{2}$ nodes.
\end{proposition}

\begin{proof}
Each pair term $q_i q_j / r_{ij}$: $\mulop(q_i, q_j) = 6\,\text{n}$ (general-domain
mul, charges can be negative), then $\divop(\cdot, r_{ij}) = 2\,\text{n}$,
with $r_{ij} > 0$.  Per pair: $6 + 2 = 8\,\text{n}$... but in v5 the scalar $k$
is free, so we focus on the addition of pair terms:
Per pair total (internal): $\approx 6\,\text{n}$.  Between-term additions
(general-domain, products of charges can be positive or negative): $2\,\text{n}$
in v5 vs $11\,\text{n}$ in v4.
Dominant saving for large $N$: $9 \binom{N}{2}$ nodes.
\end{proof}

%% =============================================================================
\section{New Algebraic Identities}
\label{sec:identities}
%% =============================================================================

\subsection{The add/sub Mirror Symmetry in $\mathcal{F}_{16}$}

\begin{theorem}[Subtraction-Addition Duality]
\label{thm:duality}
In $\mathcal{F}_{16}$, subtraction and addition are mirror constructions:
\begin{align}
  \subop(x,y) &= \LEdiv\!\bigl(x,\;\EML(y,1)\bigr) = x - \ln(e^y) = x - y, \\
  \addop(x,y) &= \LEdiv\!\bigl(x,\;\DEML(y,1)\bigr) = x - \ln(e^{-y}) = x + y.
\end{align}
Swapping $\EML$ for $\DEML$ in the inner node toggles between subtraction and
addition.  Both constructions are 2-node trees; both are valid for all
$(x, y) \in \R \times \R$.
\end{theorem}

\begin{remark}
This duality was structurally present in $\F$ from the beginning, but invisible
while $\addop$ cost $11\,\text{n}$ under v4.  The ADD-T1 result makes the
symmetry manifest: $\EML$ exponentiation with positive sign yields subtraction;
$\DEML$ exponentiation with negative sign yields addition.  Both operations
have identical node count $2\,\text{n}$, consistent with the algebraic symmetry
$x - y = x + (-y)$.
\end{remark}

\subsection{Rational Linear Combination}

\begin{theorem}[Linear combination cost]
\label{thm:linear_combo}
For rational constants $a, b \in \mathbb{Q}$:
\[
  \Cost_{\mathrm{v5}}(ax + by) = 2\,\text{n}, \qquad
  \Cost_{\mathrm{v4}}(ax + by) = 11\,\text{n} \text{ (general domain)}.
\]
More generally, for $N$ rational-coefficient terms:
$\Cost_{\mathrm{v5}}\!\bigl(\sum_{i=1}^{N} a_i x_i\bigr) = 2(N-1)\,\text{n}$.
\end{theorem}

\begin{proof}
Under the rational-terminal convention, $a \cdot x$ is a free scalar product
(leaf $x$ with rational weight; no operator node needed).
Thus $ax$ and $by$ are free.  Then:
\[
  ax + by = \LEdiv(ax,\,\DEML(by,\,1)), \quad 2\,\text{n}.
\]
For $N$ terms: $\sum_{i=1}^{N} a_i x_i$ is computed by $N-1$ sequential
$\addop$ nodes, each costing $2\,\text{n}$.  Total: $2(N-1)\,\text{n}$.
\end{proof}

\begin{corollary}[Affine neural layer, constant weights]
An affine transformation $y = \mathbf{w}^T \mathbf{x} + b$ with
$\mathbf{w} \in \mathbb{Q}^N$ costs $2(N-1) + 2 = 2N\,\text{n}$ in v5.
Under v4 add\_gen: $11N\,\text{n}$.  Saving: $9N\,\text{n}$.
\end{corollary}

\subsection{Logarithm of Sum via LEAd (Pre-existing 2-node Route)}

\begin{theorem}[$\ln(a+b)$ in 2 nodes]
\label{thm:ln_sum}
For $a > 0$, $b > -a$:
\[
  \ln(a + b) = \LEAd\!\bigl(\EXL(0,a),\;b\bigr) = \ln(e^{\ln a} + b) = \ln(a + b).
\]
Cost: $\mathbf{2\,\text{n}}$.
\end{theorem}

\begin{remark}
This 2-node route for $\ln(a+b)$ \emph{does not use} the ADD-T1 add = 2n
construction.  It predates SuperBEST v5 and uses the $\LEAd$ operator directly.
The naive route (add in $2\,\text{n}$, then $\ln$ in $1\,\text{n}$) costs $3\,\text{n}$.
Theorem~\ref{thm:ln_sum} shows the LEAd route is strictly cheaper.

Under v4 (add\_pos $= 3\,\text{n}$ for $a, b > 0$): naive route costs $3 + 1 = 4\,\text{n}$;
LEAd route costs $2\,\text{n}$.  Under v5: naive costs $2 + 1 = 3\,\text{n}$;
LEAd still costs $2\,\text{n}$.  The LEAd route wins in both versions,
but the gap narrows from $2\,\text{n}$ (v4) to $1\,\text{n}$ (v5).
\end{remark}

\subsection{Softplus and Log-Sigmoid}

\begin{proposition}[Softplus is trivially 1 node]
$\mathrm{softplus}(x) = \ln(1 + e^x) = \LEAd(x, 1)$.
$\SB(\mathrm{softplus}) = 1$.
\end{proposition}

This was always $1\,\text{n}$ in $\F$; ADD-T1 does not change it.
Equivalently: $\LSE(x, 0) = \ln(e^x + e^0) = \ln(e^x + 1) = \mathrm{softplus}(x)$.

\begin{proposition}[Log-sigmoid cost]
$\ln \sigma(x) = -\ln(1 + e^{-x})$.
\begin{align*}
  v_1 &= \DEML(x,1) = e^{-x},  & 1\,\text{n} \\
  v_2 &= \LEAd(0, v_1) = \ln(1 + e^{-x}),  & 1\,\text{n} \\
  v_3 &= \negop(v_2) = -\ln(1 + e^{-x}),  & 2\,\text{n}
\end{align*}
Total: $4\,\text{n}$.  No v5 saving (bottleneck is $\negop$, not $\addop$).
\end{proposition}

\subsection{Sum of Squares}

$x^2 + y^2$ with $x, y \in \R$ arbitrary.

\begin{proposition}
\[
  \Cost_{\mathrm{v5}}(x^2 + y^2) = 6\,\text{n}, \quad
  \Cost_{\mathrm{v4}}(x^2 + y^2) = 7\,\text{n}.
\]
\end{proposition}

\begin{proof}
$x^2 = \EML(2 \cdot \EXL(0,x), 1)$: $2\,\text{n}$ (fixed exponent).
Similarly $y^2$: $2\,\text{n}$.
$x^2 + y^2$: the two squares are positive, so add\_pos was $3\,\text{n}$ in v4
and is $2\,\text{n}$ in v5.  Total: $2 + 2 + 2 = 6\,\text{n}$ (v5);
$2 + 2 + 3 = 7\,\text{n}$ (v4).  Saving: $1\,\text{n}$.
\end{proof}

\subsection{Swish Activation}

$\mathrm{swish}(x) = x \cdot \sigma(x) = x / (1 + e^{-x})$, $x > 0$.

\begin{proposition}
\[
  \Cost_{\mathrm{v5}}(\mathrm{swish}) = 6\,\text{n}, \quad
  \Cost_{\mathrm{v4}}(\mathrm{swish}) = 7\,\text{n}.
\]
\end{proposition}

\begin{proof}
\begin{align*}
  v_1 &= \DEML(x,1) = e^{-x},  & 1\,\text{n} \\
  v_2 &= \addop(1, v_1) = 1 + e^{-x}
       = \LEdiv(1,\,\DEML(v_1,1)),  & 2\,\text{n} \;\text{(v5)} \\
  v_3 &= \ELSb(0, v_2) = 1/v_2 = \sigma(x),  & 1\,\text{n} \\
  v_4 &= \mulop(x, v_3),  & 2\,\text{n} \;(x > 0)
\end{align*}
Total v5: $1 + 2 + 1 + 2 = 6\,\text{n}$.
In v4, $\addop(1, e^{-x})$ with positive operands used add\_pos $= 3\,\text{n}$:
total $1 + 3 + 1 + 2 = 7\,\text{n}$.  Saving: $1\,\text{n}$.
\end{proof}

%% =============================================================================
\section{Geometry and Curvature}
\label{sec:geometry}
%% =============================================================================

\subsection{Gaussian Curvature}

$K = (f_{xx} f_{yy} - f_{xy}^2) / (1 + f_x^2 + f_y^2)^2$.

\begin{proposition}[Gaussian curvature cost]
\[
  \Cost_{\mathrm{v5}}(K) = 14\,\text{n}, \quad
  \Cost_{\mathrm{v4}}(K) = 16\,\text{n}.
\]
Saving: $2\,\text{n}$ (two positive-domain additions in the denominator).
\end{proposition}

\begin{proof}
\textbf{Numerator:} $f_{xx} f_{yy} - f_{xy}^2$:
$\mulop(f_{xx}, f_{yy}) = 2\,\text{n}$;
$f_{xy}^2 = 1\,\text{n}$ (fixed exponent);
$\subop = 2\,\text{n}$.  Numerator: $5\,\text{n}$.

\textbf{Denominator:} $1 + f_x^2 + f_y^2$:
$f_x^2 = 1\,\text{n}$, $f_y^2 = 1\,\text{n}$;
$\addop(1, f_x^2) = 2\,\text{n}$ in v5 (v4: $3\,\text{n}$);
$\addop(\cdot, f_y^2) = 2\,\text{n}$ in v5 (v4: $3\,\text{n}$);
$\powop(\cdot, 2) = 1\,\text{n}$.  Denominator: $1+1+2+2+1 = 7\,\text{n}$ (v5),
$1+1+3+3+1 = 9\,\text{n}$ (v4).

\textbf{Division:} $2\,\text{n}$.  Total: $5 + 7 + 2 = 14\,\text{n}$ (v5);
$5 + 9 + 2 = 16\,\text{n}$ (v4).
\end{proof}

\subsection{Von Mises Yield Criterion}

$\sigma_{VM} = \sqrt{\sigma_1^2 - \sigma_1 \sigma_2 + \sigma_2^2}$.

\begin{proposition}
\[
  \Cost_{\mathrm{v5}}(\sigma_{VM}) = 10\,\text{n}, \quad
  \Cost_{\mathrm{v4, add\_gen}}(\sigma_{VM}) = 19\,\text{n}.
\]
Saving: $9\,\text{n}$ (one general-domain addition from signed intermediate).
\end{proposition}

\begin{proof}
$\sigma_1^2 = 1\,\text{n}$; $\sigma_2^2 = 1\,\text{n}$;
$\sigma_1 \sigma_2 = 2\,\text{n}$;
$A = \subop(\sigma_1^2, \sigma_1\sigma_2) = 2\,\text{n}$
($A$ can be negative: $\sigma_1^2 - \sigma_1\sigma_2 < 0$ when $\sigma_2 > \sigma_1$);
$\addop(A, \sigma_2^2) = 2\,\text{n}$ in v5 — general-domain since $A$ can be
negative (v4: $11\,\text{n}$);
$\sqrtop = 2\,\text{n}$.
Total v5: $1+1+2+2+2+2 = 10\,\text{n}$.
Total v4: $1+1+2+2+11+2 = 19\,\text{n}$.  Saving: $9\,\text{n}$.
\end{proof}

\subsection{Mean Curvature}

$H = \bigl((1+f_y^2)f_{xx} - 2f_x f_y f_{xy} + (1+f_x^2)f_{yy}\bigr)
     / \bigl(2(1+f_x^2+f_y^2)^{3/2}\bigr)$.

\begin{proposition}[Mean curvature saving estimate]
The numerator of $H$ contains two signed general-domain additions
(between terms of mixed sign).  Each saves $9\,\text{n}$ in v5.
Total saving: $\approx 18\,\text{n}$.
\[
  \Cost_{\mathrm{v5}}(H) \approx 13\,\text{n}, \quad
  \Cost_{\mathrm{v4}}(H) \approx 31\,\text{n}.
\]
\end{proposition}

%% =============================================================================
\section{Structural Patterns}
\label{sec:patterns}
%% =============================================================================

\subsection{Elimination of the Cost-Class Split}

Under SuperBEST v4, equations were implicitly stratified into two cost classes:
\begin{itemize}
  \item \textbf{Class~P} (positive-domain sums): $\Cost = (\alpha_0 + 3)N - 3$.
  \item \textbf{Class~G} (general/signed sums): $\Cost = (\alpha_0 + 11)N - 11$.
\end{itemize}
The gap between classes was $8N$ nodes per addition-site, a factor of $11/3 \approx 3.7$.

Under SuperBEST v5:
\[
  \Cost = (\alpha_0 + 2)N - 2 \quad \text{for all sign structures.}
\]

\begin{theorem}[Unified scaling law]
\label{thm:unified_scaling}
For any $N$-term sum $\sum_{i=1}^N f(x_i)$ where each $f$-evaluation costs
$\alpha_0$ nodes in $\mathcal{F}_{16}$:
\[
  \Cost_{\mathrm{v5}} = (\alpha_0 + 2)N - 2,
\]
regardless of whether the partial sums are positive or mixed-sign.
This is a strict improvement over v4 for all equations with general-domain additions,
and a $1n$-per-addition improvement over v4 for positive-domain sums.
\end{theorem}

\subsection{New Minimum-Cost Floor for Summations}

\begin{theorem}[Cost floor for $N$-term sums]
For any $N$-term real-valued sum $\sum_{i=1}^N g_i$ where each $g_i$ has
$\mathcal{F}_{16}$-cost $\ge \alpha_0$:
\[
  \Cost \ge (\alpha_0 + 2)N - 2.
\]
This bound is \emph{tight}: achieved when additions are sequential and
each uses the ADD-T1 2-node construction.
\end{theorem}

\begin{proof}
The $N$ terms require at least $\alpha_0 N$ nodes.
The $N - 1$ additions each require at least $\SB(\addop) = 2$ nodes (ADD-T1).
Lower bound: $\alpha_0 N + 2(N-1) = (\alpha_0 + 2)N - 2$.
The ADD-T1 construction achieves equality.
\end{proof}

%% =============================================================================
\section{Summary Table}
\label{sec:summary}
%% =============================================================================

\begin{longtable}{@{}lcccc@{}}
\caption{Cost comparison: SuperBEST v4 vs v5 for key expressions.
  $k_g$ = number of general-domain add sites;
  $k_p$ = number of positive-domain add sites.
  Saving $= 9k_g + k_p$.}
\label{tab:summary}
\\
\toprule
\textbf{Expression} & \textbf{v4 cost} & \textbf{v5 cost} & \textbf{Saving} & \textbf{Type} \\
\midrule
\endfirsthead
\midrule
\textbf{Expression} & \textbf{v4 cost} & \textbf{v5 cost} & \textbf{Saving} & \textbf{Type} \\
\midrule
\endhead
\bottomrule
\endfoot
2-term LSE $\ln(e^x + e^y)$     & $6$       & $5$       & $1$       & pos \\
$N$-term LSE                    & $4N{-}3$  & $3N{-}2$  & $N{-}1$   & pos \\
Softmax ($N$ outputs)           & $6N{-}3$  & $5N{-}2$  & $N{-}1$   & pos \\
KL divergence ($N$ terms)       & $16N{-}11$& $7N{-}2$  & $9(N{-}1)$& gen \\
JSD ($N$ terms)                 & $29N{-}2$ & $17N{-}2$ & $12(N{-}1)$& gen \\
Itakura-Saito ($N$ terms)       & $18N{-}11$& $9N{-}2$  & $9(N{-}1)$& gen \\
Shannon entropy ($N$ terms)     & $8N{-}3$  & $7N{-}2$  & $N{-}1$   & pos \\
Partition function $Z$ ($N$ terms) & $6N{-}3$ & $5N{-}2$ & $N{-}1$  & pos \\
Helmholtz $F = -kT\ln Z$        & $6N$      & $5N{+}1$  & $N{-}1$   & pos \\
Gibbs $\Delta G = \Delta H - T\Delta S$ & $16$ & $7$  & $9$       & gen \\
Lennard-Jones $V_{LJ}$          & $10$      & $10$      & $0$       & sub only \\
Coulomb sum ($N$ particles)     & ${\approx}17N^2/2$ & $2N^2{-}2N{-}2$ & ${\approx}9N^2/2$ & gen \\
$x^2 + y^2$                     & $7$       & $6$       & $1$       & pos \\
$e^x + e^y$                     & $5$       & $4$       & $1$       & pos \\
$\ln(a+b)$ (LEAd route)         & $2$       & $2$       & $0$       & LEAd \\
Softplus $\ln(1+e^x)$           & $1$       & $1$       & $0$       & LEAd \\
Swish $x\sigma(x)$              & $7$       & $6$       & $1$       & pos \\
Taylor series ($N$ terms)       & $4N{-}3$  & $3N{-}2$  & $N{-}1$   & pos \\
$ax + by$ ($a,b \in \mathbb{Q}$)& $11$      & $2$       & $9$       & gen \\
$\sum_{i=1}^N a_i x_i$ (const wts) & $11N$ & $2(N{-}1)$ & $\approx 11N$ & gen \\
Von Mises $\sigma_{VM}$         & $19$      & $10$      & $9$       & gen \\
Gaussian curvature $K$          & $16$      & $14$      & $2$       & pos \\
Mean curvature $H$              & ${\approx}31$ & ${\approx}13$ & ${\approx}18$ & gen \\
DCF mixed cashflows ($N$ terms) & $20N{-}11$& $9N{-}2$  & $9(N{-}1)$& gen \\
Bilinear form ($N{\times}N$)    & $13N^2{-}11$ & $4N^2{-}2$ & $9N^2$ & gen \\
\end{longtable}

%% =============================================================================
\section{Conclusion}
\label{sec:conclusion}
%% =============================================================================

The ADD-T1 identity $x + y = \LEdiv(x, \DEML(y, 1))$ establishes
$\SB(\addop) = 2\,\text{n}$ for all real $x, y$.  The structural consequences
are far-reaching:

\begin{enumerate}
\item \textbf{Add/sub mirror symmetry.}
  Subtraction and addition are now algebraically dual 2-node constructions,
  differing only in whether the inner exponentiation is $\EML$ (positive
  exponent) or $\DEML$ (negative exponent).

\item \textbf{Unified scaling law.}
  The two-tier system (add\_pos $= 3$, add\_gen $= 11$) collapses to a single
  universal law: $\Cost(\sum_{i=1}^N f_i) = (\alpha_0 + 2)N - 2$.

\item \textbf{Rational linear combinations.}
  Any $\mathbb{Q}$-linear combination $\sum a_i x_i$ costs only $2(N-1)$
  nodes in v5, down from $11(N-1)$ in v4.  This dramatically lowers the
  floor for linear algebra in $\F$.

\item \textbf{Pre-existing 2n route for $\ln(a+b)$.}
  The $\LEAd$ operator gives $\ln(a+b) = \LEAd(\ln a, b)$ in $2\,\text{n}$,
  beating the naive add-then-ln route ($3\,\text{n}$ in v5, $4\,\text{n}$ in v4).

\item \textbf{Large-scale beneficiaries.}
  Expressions with many general-domain additions (KL divergence, Coulomb sums,
  bilinear forms, mean curvature) save $9\,\text{n}$ per addition site.
  For $N$-term expressions this is a saving of $\Theta(N)$ or $\Theta(N^2)$ nodes.
\end{enumerate}

\noindent
The SuperBEST v5 table is complete and optimal.  The next frontier is whether
extending $\mathcal{F}_{16}$ to $\mathcal{F}_{17}$ or higher could reduce
any entry below its current bound.

\bigskip
\noindent\textit{Almaguer, A.R.\ (2026).
New Algebraic Identities and Minimal Constructions Under add $= 2$ Nodes.
monogate research.\ \url{https://monogate.org}}

\begin{thebibliography}{9}

\bibitem{ADDT1}
Almaguer, A.R.\ (2026).
\textit{General Addition in Two Nodes: Closing the Last Gap in SuperBEST v4}.
\texttt{python/paper/theorems/ADD\_T1\_General\_Addition\_2n.tex}.

\bibitem{v5Audit}
Almaguer, A.R.\ (2026).
\textit{SuperBEST v5: Complete Structural Proof of Optimality for All Ten Core
Arithmetic Operations}.
\texttt{python/paper/theorems/SuperBEST\_v5\_Structural\_Audit.tex}.

\bibitem{CascadeV5}
Almaguer, A.R.\ (2026).
\textit{Cascade Effects of add $= 2n$ on the SuperBEST Cost Catalog}.
\texttt{python/paper/cascade\_v5\_ripple.tex}.

\bibitem{ScalingLaws}
Almaguer, A.R.\ (2026).
\textit{Scaling Laws Under SuperBEST v5}.
\texttt{python/results/scaling\_laws\_v5.json}.

\bibitem{ELCBoundary}
Almaguer, A.R.\ (2026).
\textit{ELC(R) Boundary Map: Definitive Three-Tier Classification}.
\texttt{python/results/x15\_elc\_boundary\_map.json}.

\end{thebibliography}

\end{document}
